%
% Documento: Citações
%

\section{Citações (ABNT NBR 10520)}

Esta seção apresenta exemplos práticos de citação conforme a ABNT NBR 10520.

\subsection{Citação direta curta (até 3 linhas)}
Deve permanecer no corpo do parágrafo, entre aspas duplas, com indicação de autor, ano e página:
\begin{flushleft}
\textit{Ex.:} De acordo com \citeonline[p.~249]{marshall1980}, ``não se pode discutir se o valor e os preços são governados pela utilidade ou pelo custo de produção''.
\end{flushleft}

Outras variações úteis:
\begin{itemize}
  \item Autor no texto + ano entre parênteses: \citeauthor{marshall1980} (\citeyear{marshall1980}) afirma que ``\ldots'' \cite[p.~249]{marshall1980}.
  \item Apenas referência entre parênteses: ``\ldots'' \cite[p.~249]{marshall1980}.
\end{itemize}

\subsection{Citação direta longa (mais de 3 linhas)}
Deve ir em bloco destacado, com recuo de 4\,cm da margem esquerda, fonte menor e sem aspas. 
No \texttt{abntex2}, use o ambiente \texttt{citacao}, que já aplica o padrão:
\begin{citacao}
Da mesma forma que não se pode afirmar se é a lâmina inferior ou superior de uma tesoura que corta uma folha de papel, também não se pode discutir se o valor e os preços são governados pela utilidade ou pelo custo de produção. (MARSHALL, 1980, p.~249)
\end{citacao}

\subsection{Citação indireta (paráfrase)}
Reproduz a ideia do autor com suas palavras. Página é opcional; recomenda-se quando ajuda o leitor.
\begin{flushleft}
\textit{Ex.:} A produção acadêmica sobre varejo no Brasil é desproporcional à relevância do segmento na economia \cite{angelo1993}.
\end{flushleft}
Variações:
\begin{itemize}
  \item \citeauthor{angelo1993} (\citeyear{angelo1993}) argumentam que \ldots
  \item Conforme apontado na literatura \cite{tanenbaum1995sistemas,de2017sistema}, \ldots
\end{itemize}

\subsection{Citação de citação (apud)}
Use somente quando o original não puder ser consultado. Indique o autor original, seguido de ``apud'' e do autor efetivamente consultado, com página.
\begin{flushleft}
\textit{Ex. no corpo do texto:} Enguita (apud \citeauthor{silva1991}, \citeyear{silva1991}, p.~21) chegou às mesmas conclusões.
\end{flushleft}

Também é possível registrar em nota de rodapé a fonte consultada:
\begin{flushleft}
\textit{Ex. com nota:} Enguita\footnote{Informação citada por \citeauthor{silva1991} (\citeyear{silva1991}, p.~21).} chegou às mesmas conclusões.
\end{flushleft}

\subsection{Comandos úteis (abntex2cite)}
\begin{itemize}
  \item \verb|\cite{chave}| → (AUTOR, ano)
  \item \verb|\citeonline{chave}| → AUTOR (ano)
  \item \verb|\citeauthor{chave}| → AUTOR
  \item \verb|\citeyear{chave}| → ano
  \item \verb|\cite[p.~10]{chave}| → (AUTOR, ano, p.~10)
\end{itemize}

\subsection{Observações finais}
\begin{itemize}
  \item Para siglas: escreva por extenso na primeira ocorrência, seguidas da sigla: Associação Brasileira do Trabalhador (ABT).
  \item As referências completas devem aparecer na lista de referências ao final do trabalho.
\end{itemize}

